\section{目的，保証，読み方}

\subsection{解く問題}

Egison coreではmatcherがfirst-class valueであり，式はmatcherを生成することも，matcherを必要とする
位置へ渡されることもある．この二つを同じ型だけで表すと，通常の型等価性と暗黙coercionのどちらを
選ぶかが，unificationの失敗や特定のsource構文に依存してしまう．本仕様はproducer type
$\Matcher{\kappa}{\tau}$とconsumer expectation
$\MatcherSlot{\kappa}{\tau}$を分け，resolved expected headがslotであるchecking cutだけを
coercion demandとする．

このcalculusが同時に満たすべき目標は次の四つである．

\begin{enumerate}
\item \textbf{source acceptanceを一つにする．}
  programが型付くことは$\mathsf{SourceTyping}$だけで定め，動的メタ理論用の内部invariantを
  第二のsource type systemにしない．
\item \textbf{推論を実行可能かつfail closedにする．}
  raw traversalは停止し，有限validatorを通った結果だけを公開する．callerにsolver traceや
  bridge certificateを要求しない．
\item \textbf{coercionをdemand-directedかつno-guessにする．}
  未解決型をcoercionのためにslotへ推測せず，通常単一化の失敗後に別branchを試さない．
\item \textbf{source typingを動的安全性へ接続する．}
  推論順序に依存する状態を消去し，value typing，preservation，matching safetyが消費できる
  state-freeな$\mathsf{TypingInvariant}$を構成する．
\end{enumerate}

\subsection{公開保証}

公開推論器とsource typingの二方向，およびclosed programの安全性は次の形で機械化されている．
\[
\begin{array}{rcl}
\mathsf{infer}(\Sighat,\Gamma,e)=\mathsf{some}\;r
&\Longrightarrow&
\mathsf{SourceTyping}(\Sighat,\Gamma,e,r.\mathsf{resolvedTarget}),\\[0.3em]
\mathsf{SourceTyping}(\Sighat,\Gamma,e,\tau)
+\mathsf{FrozenSigWF}(\Sighat)
&\Longrightarrow&
\mathsf{infer}(\Sighat,\Gamma,e)\ne\mathsf{none},\\[0.3em]
\mathsf{SourceTyping}(\Sighat,\emptyset,e,\tau)
+\mathsf{FrozenSigWF}(\Sighat)
&\Longrightarrow&
\mathsf{SafeResult}_{\mathsf{Source}}(\Sighat,e,\tau,\mathcal S_F).
\end{array}
\]
従って$\mathsf{FrozenSigWF}$の下で，一般contextのsource typabilityと公開推論の成功は同値である．
ただし$\mathsf{SourceTyping}$自体がterminal auditを定義に含む．従ってこの同値と受理完全性は，audit済みの
$\mathsf{SourceTyping}$に相対的であり，raw derivation，auditを除いた独立の宣言的型付け，または実行時に
安全な全programに対する完全性ではない．
同じsourceの二つの$\mathsf{SourceTyping}$ targetは，全residual capability／target metavariableの
局所renamingを法として一意である．さらに，source targetに現れる二sort metavariableだけを変更する
有限support substitutionでinstance preorderを定めると，$\mathsf{inferType}$の返値は一般contextでも
principalである．open termについては，実行runのresolved context／targetと，各derivationの
terminal-normalized context／targetを同じpaired substitutionで比較する相対principalityも成立する．

pattern-freeでcapability-inertなDamas--Milner断片についても二方向の境界を機械化する．closed
source fragmentのtypingはDM typingへ消去でき，逆に任意のDM typing derivationは，そのcontextを
埋め込んだ公開推論器で受理される：
\[
\mathsf{FrozenSigWF}(\Sighat)\land
\mathsf{DM.Typing}(\Gamma,e,\tau)
\Longrightarrow
\mathsf{inferenceSucceeds}(\Sighat,\mathsf{emb}(\Gamma),e)=\mathsf{true}.
\]
この保証はacceptanceであり，DM derivationが選ぶ$\tau$と公開推論器の返値型の構文的一致ではない．
後者はprincipality／instance preorderによる別の比較問題である．

\subsection{判断の役割を分ける}

\begin{center}
\begin{tabular}{p{0.25\linewidth}p{0.65\linewidth}}
\toprule
構成要素 & 役割 \\
\midrule
$\mathsf{SourceTyping}$
  & terminal auditを含み，source acceptanceを定義する唯一の公開judgment \\
$\mathsf{ExprDeriv}$
  & 成功した実行traceを規則木へ戻すproof-relevantな内部reconstruction \\
$\mathsf{TypingInvariant}$
  & supply，substitution，ledgerを消去した動的メタ理論用の内部invariant \\
$\mathsf{ValueTy}$とmatching judgments
  & evaluationとmatching machineのpreservation／progressを述べるruntime側の判断 \\
\bottomrule
\end{tabular}
\end{center}

依存方向は
\[
\mathsf{SourceTyping}
\xRightarrow{\mathsf{state\ erasure}}
\mathsf{TypingInvariant}
\xRightarrow{\mathsf{dynamic\ metatheory}}
\mathsf{runtime\ safety}
\]
である．$\mathsf{TypingInvariant}$の存在から$\mathsf{SourceTyping}$や推論成功を逆向きに推論しない．

\subsection{非目標と現在の境界}

一般のprogram terminationは主張しない．recursionはsingleton direct-selfの単相$\mathtt{fix}$に
限る．matcher literalはshape，catch-all order，binder線形性，data-arm exhaustiveness，coverageを
すべて要求する．未解決lambda domainを複数箇所で共有したときのsource-order依存は正負回帰で
固定する．これはno-guessとchronological state threadingの意図された言語境界であり，source要素の
順序不変性は主張しない．
またDM acceptanceは推論成功だけを主張し，DM targetと推論返値の構文的一致は非目標である．

terminal auditには，rawでは受理されるが公開推論では拒否されるという，意図した言語境界とは区別すべき
既知の受理損失がある．整形式なsignatureの下で，
一要素tupleの要素を引数matcherへ委譲する閉じたlambdaはraw推論に成功するが，引数slotから借りた
capability variableが局所finalization後に観察不能な$\mathsf K$へ特殊化される．9個のterminal checkは
\[
(\mathtt{true},\mathtt{true},\mathtt{true},\mathtt{true},\mathtt{true},
 \mathtt{true},\mathtt{true},\mathtt{false},\mathtt{true})
\]
であり，matcher-finalization checkだけが失敗して公開推論は拒否する．compile-time guardがこの具体計算を
検査し，通常の定理が任意のraw成功等式からdemand-directed derivationを構成する．閉lambdaがclosureへ
評価されることも証明済みであるが，auditと独立な宣言的型付けに基づく安全性は未証明である．従ってこの例を
一般的な意味で安全とする定理や，terminal auditによる受理損失の正確な範囲は公開保証に含めない．

修正候補は，matcher自身が観察した構造と，引数slotから完全な値として借りたleafの由来をfinalizationまで
保持することである．借りたvariableも一律にfreezeする方法は意図した後続特殊化を失い，すべてのopaque
leafを一律に許す方法はmatcher自身が内部構造を観察する場合まで受理し得るため，どちらも採用しない．

\subsection{読み方}

第2章は，二sort型，三つの推論状態，Origin certificate，terminal audit，validatorが必要になる
理由と，公開定理までの依存関係を示す．第3章で構文と判断を定義し，第4章で一回のchecking cutを
定式化する．第5章は通常の式，第6章はpatternとmatcher literalの規則である．第7章は局所代数を
traversalへ持ち上げ，soundness，acceptance completeness，state erasure，安全性，受理同値，target
一意性，principalityへ接続する．Lean moduleと回帰の索引は\texttt{../docs/details.md}に譲る．

設計だけを追う読者は第2章，第4章，第7章を読めばよい．規則を実装と照合する場合は第3--6章，
証明の依存を確認する場合は第2.4節と第7章を読む．
